\documentclass[10pt,a4paper]{amsart}
\usepackage[margin=19mm]{geometry}
\usepackage{amsmath,amssymb,mathtools}
\usepackage[scr=rsfs,cal=euler]{mathalfa}
\usepackage{xfrac}
\usepackage[unicode,hidelinks,bookmarks=false]{hyperref}
\newcommand{\ie}{i.e.\ }
\newcommand{\cf}{cf.\ }
\newcommand{\resp}{respectively\ }
\newcommand{\Subsection}{\subsection}
\providecommand{\nobreakdash}{\nobreak\hbox{-}}
\setlength{\emergencystretch}{2em}
\multlinegap0pt
\usepackage{amssymb}
\usepackage{mathtools}
\usepackage{bm}
\usepackage{algorithm}\usepackage{algpseudocode}
\newtheorem{proposition}{Proposition}
\newcommand{\bL}{\bm{L}}
\newcommand{\bu}{\bm{u}}
\newcommand{\bw}{\bm{w}}
\newcommand{\cU}{\mathcal{U}}
\title{End-to-End Learning of Correlated Operating Reserve Requirements in Security-Constrained Economic Dispatch}
\author{Owen Shen, Hung-po Chao, Haihao Lu, Patrick Jaillet}
\date{Source: arXiv:2604.05167v1 (CC BY 4.0)}
\begin{document}
\maketitle

\noindent\emph{Source and licence.} End-to-End Learning of Correlated Operating Reserve Requirements in Security-Constrained Economic Dispatch. Version arXiv:2604.05167v1 (CC BY 4.0), distributed by arXiv under the Creative Commons Attribution 4.0 International licence, http://creativecommons.org/licenses/by/4.0/, as recorded in the arXiv licence metadata for this exact version and shown on the abstract page https://arxiv.org/abs/2604.05167v1. Full licence notice: \texttt{licenses/arxiv-2604.05167-CC-BY-4.0.txt}.\par\smallskip

\noindent\textit{Editorial scope.} Counted unit: the article's proposition prop:ellipsoid-grads (source0.tex lines 786--791). The article's complete original proof of it is delivered with it (source0.tex lines 793--797, one \texttt{proof} environment of 5 lines). No mathematics is written, repaired or re-derived: the statement and the proof are the article's own lines, in the article's own notation, and no retained line is substituted or re-flowed.  No further source-local context is needed: every cross-reference of every retained line resolves inside the two delivered spans.  Nothing was omitted from any retained span, so the package carries no omission record.  No retained line contains a citation, so the wrapper carries no bibliography and no bibliography entry.  No retained line contains a figure environment, an image-inclusion command or drawing code, so no image is delivered and the package declares no asset dependency.  Editorial formatting only, in addition: an AMS \texttt{amsart} preamble that restates the article's own declarations and macros used by the retained lines, namely the environments \texttt{proposition} on separate counters, together with the standard packages those lines need.  The retained environments print in this document as Proposition 1 in order of appearance, whereas the article prints the same items with the numbering of its own full document.  The complete native source of the article as distributed in the arXiv e-print for this exact version is bundled unchanged as \texttt{sources/197934/source0.tex}.
\begin{proposition}[Ellipsoidal Gradients]\label{prop:ellipsoid-grads}
For $\bw\neq 0$ with $\bL^\top\bw\neq 0$:
$\nabla_{\bL}\sigma_{\cU_{\bL,\rho}}(\bw) = \rho\bw\bw^\top\bL/\|\bL^\top\bw\|_2$.
For $\bu\neq 0$:
$\nabla_{\bL}s_{\bL}(\bu) = -\bL^{-\top}(\bL^{-1}\bu)(\bL^{-1}\bu)^\top/\|\bL^{-1}\bu\|_2$.
\end{proposition}

\begin{proof}
For the support function, let $y:=\bL^\top\bw$. Then $\mathrm{d}\sigma = \rho y^\top\mathrm{d}y/\|y\|_2 = \rho(\bL^\top\bw)^\top(\mathrm{d}\bL)^\top\bw/\|\bL^\top\bw\|_2$, which identifies the Frobenius gradient as $\rho\bw(\bL^\top\bw)^\top/\|\bL^\top\bw\|_2 = \rho\bw\bw^\top\bL/\|\bL^\top\bw\|_2$.

For the gauge, let $v:=\bL^{-1}\bu$. Then $\mathrm{d}v=-\bL^{-1}(\mathrm{d}\bL)v$ gives $\mathrm{d}s = -v^\top\bL^{-1}(\mathrm{d}\bL)v/\|v\|_2$, identifying the gradient as $-\bL^{-\top}vv^\top/\|v\|_2$.
\end{proof}

\end{document}
